
%============================================================
\part{Experimental Database}
\label{part:experimental-database}
%============================================================


%============================================================
%\chapter{Experimental Database: Scientific Objectives and Principles}
%\label{chap:experimental-objectives}
%============================================================

\chapter[Experimental Database: Scientific Objectives and Principles]
{Experimental Database: Scientific Objectives and Principles}
\chaptermark{Experimental Database}
\label{chap:experimental-objectives}


\section{Introduction and Scope}

The CIVITA Experimental Database is a community-driven scientific
resource dedicated to the systematic collection, documentation,
curation, preservation, and dissemination of experimental data
on concrete and cementitious materials. Its primary purpose is
to establish a comprehensive and reliable body of experimental
evidence that can support the development, calibration, validation,
and objective assessment of computational models.

Experimental investigations of concrete materials and structures
have generated an extensive body of scientific information over
several decades. However, much of this information remains
distributed across journal publications, technical reports,
dissertations, institutional repositories, and laboratory archives.
Experimental results are frequently reported using heterogeneous
terminology, incomplete metadata, different measurement conventions,
and inconsistent formats. In many cases, the information required
to reproduce an experiment computationally is difficult to retrieve
or reconstruct.

These limitations significantly reduce the potential value of
existing experimental evidence. Even when numerical results are
available, their interpretation often requires additional
information concerning material composition, specimen preparation,
geometry, reinforcement configuration, environmental conditions,
boundary conditions, loading histories, and measurement procedures.
The absence of such information may prevent meaningful comparisons
between experimental observations and computational predictions.

The CIVITA Experimental Database addresses these challenges through
a standardized framework for representing experimental information
and the relationships among its constituent elements. The database
is intended to accommodate investigations conducted at different
spatial and structural scales, from material characterization
to the experimental assessment of structural components and
systems. Its scope includes mechanical, physical, thermal,
hygral, chemical, and other phenomena relevant to the behavior
and durability of concrete and cementitious materials.

The database accommodates both newly generated experimental
information and existing datasets recovered from the scientific
literature. Experimental observations may include scalar
measurements, curves, time histories, tabular data, images,
videos, and multidimensional fields. Each dataset is accompanied
by structured metadata providing the scientific context
necessary for its interpretation and reuse.

The database architecture and scientific conventions are
defined independently of the applications used to contribute,
curate, or retrieve experimental information. The corresponding
data model, representation standards, submission procedures,
curation processes, and search capabilities are documented
in the subsequent chapters of this Part.


\section{Scientific Objectives}

The principal scientific objective of the CIVITA Experimental
Database is to transform distributed and heterogeneous experimental
information into a structured, traceable, and reusable collection
of scientific evidence.

This objective is pursued through several complementary goals.

First, the database seeks to preserve experimental observations
together with sufficient contextual information to support
their scientific interpretation. Material characteristics,
specimen configurations, experimental procedures, environmental
conditions, and measured responses are represented as related
scientific entities rather than isolated data records. This
organization facilitates the reconstruction of experimental
conditions and the preparation of corresponding computational
simulations.

Second, the database promotes consistency in the representation
of experimental information through shared scientific definitions,
controlled vocabularies, standardized parametrizations, explicit
units, and machine-readable metadata. These conventions reduce
ambiguity and enable meaningful comparisons among datasets
originating from different laboratories, experimental programs,
and reporting practices.

Third, the database aims to facilitate the discovery and
systematic selection of experimental evidence. Researchers
should be able to identify datasets satisfying specific
requirements concerning materials, specimen dimensions,
reinforcement characteristics, experimental conditions,
observed phenomena, and measured quantities. Such capabilities
are essential for assembling experimental datasets suitable
for calibration, validation, and comparative model assessment.

Fourth, the database provides the experimental foundation
for reproducible computational modeling workflows. Experimental
observations and their associated metadata should be accessible
through structured representations that can be interpreted
and processed by computational tools with minimal ambiguity.
This includes the ability to connect experimental records
with model parameter identification procedures, independent
validation exercises, and other CIVITA scientific objects.

Finally, the database seeks to preserve experimental knowledge
over time and facilitate its cumulative expansion through
community contributions. This includes recovering valuable
historical datasets, documenting newly generated experimental
evidence, correcting or extending existing records, and
progressively improving the completeness and quality of
available scientific information.


\section{Guiding Principles}

The design and development of the CIVITA Experimental Database
are guided by scientific principles intended to ensure rigor,
transparency, interoperability, and long-term usability.
The database is conceived in accordance with the FAIR data
principles, according to which scientific information should
be Findable, Accessible, Interoperable, and Reusable.

These principles are implemented through the following
interrelated requirements.

\paragraph{Findability.}

Experimental information is organized through standardized
identifiers, structured metadata, controlled vocabularies,
and explicit relationships among scientific entities.
These features support indexing, parameter-based searches,
and the identification of experimental records satisfying
specified scientific criteria. Findability extends to
individual experimental tests and associated observations,
allowing users to discover relevant evidence within larger
experimental programs.

\paragraph{Accessibility.}

The database is intended to provide systematic access to
experimental metadata and associated scientific data,
subject to applicable access and redistribution conditions.
Structured records and references to supporting files
facilitate retrieval and interpretation without requiring
users to reconstruct their relationships manually.
The adopted representations are designed to remain
independent of any particular user interface or
storage technology.

\paragraph{Interoperability.}

Experimental information is represented through consistent
scientific definitions, explicit units, standardized
parametrization schemes, controlled vocabularies, and
machine-readable schemas. These conventions enable data
exchange among different CIVITA components and facilitate
integration with external computational and analytical tools.

\paragraph{Reusability.}

Experimental datasets are documented with sufficient
information to support scientific applications beyond
those originally considered by their authors. Reusability
requires preservation of experimental context, clear
definitions of reported quantities, explicit provenance,
and consistent representation of data and metadata.
Particular emphasis is placed on supporting computational
model development, parameter identification, validation,
uncertainty quantification, and comparative assessments.

\paragraph{Scientific traceability.}

Experimental information must remain traceable to its
original source and, where applicable, to the procedures
used for extraction, interpretation, transformation,
and curation. This requirement is especially important
when experimental observations are recovered from
published figures or tables rather than obtained
directly from original laboratory records.

\paragraph{Extensibility.}

The database is designed to accommodate new materials,
experimental methods, physical phenomena, measurement
technologies, and data representations. Its scientific
definitions and controlled vocabularies may evolve
through documented extensions and revisions while
preserving the interpretation of previously registered
scientific information.

\paragraph{Separation of scientific definitions and data.}

The scientific meaning of experimental entities,
parameters, and relationships is established through
shared specifications independently of the applications
and technologies used to store or process individual
datasets. Structured metadata describe experimental
information and provide references to associated
scientific data, allowing the database infrastructure
to evolve without altering the underlying scientific
definitions.


\section{Data Provenance, Traceability, and Quality}

The scientific value of experimental information depends
on the reliability of the observations and on the
completeness and accuracy of the documentation describing
the conditions under which they were obtained.

The CIVITA Experimental Database therefore places particular
emphasis on preserving the provenance of experimental
data throughout their acquisition, extraction, processing,
curation, and subsequent reuse.

Each contributed dataset is associated with an identifiable
scientific source, which may correspond to a publication,
technical report, dissertation, experimental program,
or another documented source of experimental evidence.
Where available, bibliographic information, digital
identifiers, and references to original figures, tables,
or supplementary materials should be preserved.

A distinction is maintained between experimental information
reported directly by the original investigators and
information subsequently extracted, reconstructed,
processed, or interpreted by data contributors. For
example, numerical values obtained by digitizing a
published load--displacement curve may differ from
the original measurements because of graphical resolution,
digitization procedures, and other sources of uncertainty.
The method used to obtain such values should therefore
be documented whenever possible.

Similarly, derived quantities should be distinguishable
from directly measured quantities, and any transformations
or assumptions introduced during data processing should
remain identifiable. This distinction is essential for
avoiding the attribution of uncertainties associated with
data reconstruction to the underlying experiment.

Data quality is considered through several complementary
dimensions, including completeness, consistency, traceability,
and scientific plausibility. Completeness concerns the
availability of information required to interpret an
experimental observation. Consistency concerns agreement
among related records, units, identifiers, and reported
quantities. Traceability concerns the ability to recover
the origin and processing history of the information.
Scientific plausibility concerns the compatibility of
the reported data with the documented experimental
configuration and established physical constraints.

The database accommodates datasets with different levels
of completeness, particularly when recovering historical
experimental results for which some information is no
longer available. Missing information should be explicitly
distinguished from quantities that are not applicable
or were not measured. Unreported values should not
be replaced by assumptions or estimates without
appropriate documentation.

Automated validation can identify structural inconsistencies,
missing required fields, invalid parameter values,
and violations of prescribed relationships. However,
compliance with a machine-readable schema does not
establish the scientific correctness of experimental data.
Scientific review remains necessary to assess the
credibility of extracted information and the suitability
of the associated documentation.

The procedures governing automated validation, scientific
review, correction, and publication are discussed in
Chapter~\ref{chap:experimental-curation}.


\section{Community Contributions and Scientific Reuse}

The CIVITA Experimental Database is conceived as a
collaborative scientific resource whose long-term
value depends on sustained contributions from the
concrete research community.

Contributions may originate from researchers providing
data generated through ongoing experimental programs,
from investigators sharing previously collected
laboratory datasets, or from contributors recovering
experimental information reported in existing scientific
publications.

The systematic recovery of historical experimental
evidence is particularly important. A substantial
fraction of experimental knowledge relevant to concrete
mechanics and structural behavior is available only
through published figures, tables, and descriptive
information. When adequately documented, these
observations can provide valuable evidence for the
assessment of computational models, including
phenomena and experimental configurations that
may be difficult or costly to reproduce.

The database adopts a common scientific representation
independently of the method used to prepare a contribution.
Data submission procedures may involve direct manual
entry, manual extraction and digitization, or
AI-assisted extraction and preparation. All modalities
are intended to produce information conforming to
the same scientific specifications and validation
requirements.

AI-assisted procedures may facilitate the recovery
of information from scientific publications and the
organization of complex experimental records. Their
use nevertheless requires careful verification,
particularly when extracting numerical measurements,
interpreting experimental configurations, or identifying
relationships among scientific entities. The provenance
of extracted data and the responsibility for reviewing
their scientific accuracy remain essential regardless
of the submission modality.

The database also supports community-driven development
of its scientific specifications. Contributors may
identify limitations in existing definitions or
propose additional material descriptors, experimental
configurations, parametrization schemes, controlled
vocabulary entries, and measurement types.

Such extensions should be evaluated for scientific
consistency and compatibility with existing definitions
before their incorporation into the authoritative
specifications.

Scientific reuse is encouraged through the preservation
of source attribution, explicit metadata, and
machine-readable relationships among experimental
entities. Researchers using CIVITA datasets remain
responsible for evaluating their suitability for
the intended scientific application and for
acknowledging the original sources and applicable
contributions.

The data submission procedures are described in
Chapter~\ref{chap:experimental-submission}, while
the search and retrieval capabilities are discussed
in Chapter~\ref{chap:experimental-search}.


\section{Relationship with Other CIVITA Components}

The Experimental Database constitutes the principal
repository of physical evidence within the CIVITA
Network. Its scientific value is enhanced by its
interoperability with the Computational Database
and the Parameter Identification, Validation, and
Benchmarking modules.

Within the Parameter Identification environment,
experimental datasets provide the observations
used to determine model parameters. The structured
representation of experimental configurations,
measured quantities, and associated conditions
facilitates the definition of calibration problems
and the comparison between observed and simulated
responses.

Within the Validation environment, experimental
data provide the independent evidence required
to assess the predictive capability of previously
calibrated computational models. Traceable experimental
records make it possible to document the conditions
under which comparisons are performed and to
distinguish uncertainties associated with experimental
observations from discrepancies attributable to
computational predictions.

The Computational Database provides the complementary
representation of the mathematical models, numerical
implementations, parameter configurations, and
simulation results used to interpret experimental
evidence. Interoperability between the two databases
supports consistent associations between experimental
configurations and their computational counterparts.

The Benchmarking environment is primarily concerned
with standardized computational problems and the
comparison of numerical results. Experimental evidence
may contribute to the scientific motivation and
definition of selected benchmarks, although
benchmarking exercises do not necessarily require
direct comparison with experimental observations.

The integration of these components is intended to
support traceable scientific workflows in which
experimental information can be selected, transferred,
interpreted, and reused without repeatedly reconstructing
its underlying scientific context.

By providing structured and persistent experimental
evidence, the CIVITA Experimental Database establishes
a foundation for cumulative model assessment and
the progressive development of reliable predictive
tools for concrete materials and structures.


%============================================================
\chapter{Experimental Database: Data Model and Hierarchy}
\label{chap:experimental-data-model}
\chaptermark{Experimental Database: Data Model}
%============================================================

\section{Overall Database Organization}

The CIVITA Experimental Database adopts a hierarchical data
model designed to represent experimental investigations
through interconnected scientific entities. The hierarchy
preserves the relationships among experimental sources,
material systems, specimen configurations, testing
procedures, and measured responses, while minimizing
redundancy and supporting the systematic reuse of
scientific information.

The underlying principle is the decomposition of an
experimental investigation into logically distinct
components, each representing a specific aspect of
the experimental program. For example, a single
material batch may be used to manufacture several
specimens with different geometries, while specimens
with identical configurations may be subjected to
different experimental tests. Similarly, a single
test may produce multiple experimental observations,
including scalar quantities, load--displacement curves,
time histories, and full-field measurements.

Representing these elements as separate but related
entities makes it possible to preserve their scientific
meaning and reconstruct their relationships without
repeatedly storing identical information.

The database hierarchy is illustrated in
Fig.~\ref{fig:database_hierarchy}.

\begin{figure}[htbp]
    \centering
    \includegraphics[width=0.9\linewidth]
    {./figs/database_hierarchy.png}
    \caption{Hierarchical organization of the CIVITA
    Experimental Database, showing the relationships
    among scientific sources, material batches,
    geometries, specimens, reinforcement systems,
    experimental tests, devices and supports,
    loading histories, and experimental observations.}
    \label{fig:database_hierarchy}
\end{figure}

The current data model comprises thirteen entity
collections, organized around a common literature
source and connected through explicitly defined
relationships. Each collection contains records
representing entities of the corresponding scientific
category.

The hierarchy distinguishes three principal
organizational levels.

The first level establishes the scientific source
and the reusable definitions associated with
the experimental investigation. These include
material batches, geometries, reinforcement systems,
and additive manufacturing variables.

The second level represents specimen configurations
and experimental tests. Specimens associate reusable
material and geometric definitions with specific
experimental configurations, while tests describe
the experimental procedures and conditions applied
to those specimens.

The third level contains information associated
with individual tests, including devices and
supports, loading histories, and experimental
observations.

This organization supports the representation of
experimental programs involving multiple material
mixtures, specimen configurations, testing conditions,
and measured responses within a consistent data
structure.

The scientific definitions and relationships
described in this chapter are independent of
the applications used for data submission,
curation, storage, or retrieval. The corresponding
machine-readable specifications are maintained
in the CIVITA Experimental Database Schema repository,
while detailed field definitions are documented
in Appendix~\ref{app:experimental-schema}.


\section{Scientific Entities and Their Relationships}

The database hierarchy is based on the distinction
between scientific entities, their attributes,
and the relationships connecting them.

A scientific entity represents an identifiable
component of an experimental investigation,
such as a material batch, specimen configuration,
experimental test, or measured response.
Its attributes describe the corresponding
physical, geometric, procedural, or observational
characteristics.

Relationships establish how entities are
associated within an experimental program.
These relationships are represented through
identifiers that connect individual records,
allowing the database to preserve the logical
organization of the experiment.

A fundamental distinction is maintained between
reusable definitions and entities associated
with individual experimental procedures.

For example, a geometry record describes
a specimen configuration that may be referenced
by several specimens. Similarly, a reinforcement
material may be associated with multiple
reinforcement components, and a material batch
may be used for several specimen configurations.

This approach avoids unnecessary duplication
and ensures that information shared among
multiple experimental records remains consistent.

The hierarchical organization is complemented
by relationships that connect different branches
of the data model. For example, reinforcement
layouts are associated with specimen geometries,
reinforcement components reference reinforcement
materials, and experimental data may be
associated with specific measurement devices.

These relationships are essential for preserving
the scientific interpretation of experimental
information. A measured displacement, for example,
may require identification of the specimen,
the experimental test, the loading configuration,
and the device or measurement procedure
associated with the observation.

The database also distinguishes the definition
of an experimental configuration from the
individual observations obtained from its
repeated execution. This distinction permits
multiple experimental repetitions to be
represented without unnecessarily duplicating
their common configuration.

The individual entity collections and their
principal relationships are described in the
following sections.


\section{Literature Sources and Keywords}

The \texttt{literature\_sources} collection
constitutes the highest organizational level
of the Experimental Database.

Each literature source represents an identifiable
origin of experimental information, such as a
journal publication, conference paper, technical
report, dissertation, published dataset, or
documented experimental campaign.

A source record establishes the scientific
provenance of the associated information and
provides the common reference through which
the corresponding material, specimen, test,
and measurement records are organized.

Relevant metadata may include bibliographic
information, authorship, publication details,
digital identifiers, and other descriptors
necessary to identify the original source.

The source establishes the scope of a dataset
submission. Under the current database hierarchy,
each submission contains exactly one source record,
and all associated entity records belong to
that same source.

A single source may nevertheless contain
multiple material batches, specimen configurations,
experimental tests, and associated observations.

The \texttt{keywords} collection provides
searchable descriptors associated with the
source or dataset. Keywords facilitate the
identification and classification of experimental
information according to relevant scientific
subjects, materials, phenomena, testing methods,
and other characteristics.

Unlike most other entity collections, individual
keyword records are identified by their keyword
values rather than by a numerical identifier
with a dedicated prefix.

Source-level metadata and keywords complement
the more detailed scientific descriptors
contained in the associated entity records.
Their role is to support discovery and
attribution without replacing the structured
representation of the underlying experimental
information.


\section{Material Batches}

The \texttt{material\_batches} collection
describes the material systems investigated
within an experimental source.

A material batch represents a particular
concrete mixture, cementitious composition,
or other material system associated with
the preparation and testing of specimens.

Material batch records provide the principal
connection between experimental observations
and the composition and characteristics of
the corresponding materials.

Depending on the material system and the
information available, the description may
include mixture proportions, constituent
materials, aggregate characteristics,
cementitious components, admixtures,
fibers, and relevant preparation or
processing conditions.

Additional information may characterize
the material state, preparation history,
curing conditions, and other attributes
necessary for interpreting experimental
responses.

A single experimental source may include
multiple material batches representing
different mixture compositions or material
systems. Each batch is assigned a distinct
identifier within the corresponding source.

A material batch may be referenced by
multiple specimen records. Consequently,
the same material definition can be
associated with specimens having different
geometries, reinforcement configurations,
or manufacturing characteristics.

This organization separates material
characterization from specimen geometry
and testing conditions, facilitating
comparisons of experimental responses
associated with changes in material
composition while preserving the
corresponding experimental context.

The material batch definition may include
information originally reported by the
investigators as well as documented
information recovered during subsequent
data extraction or curation. The distinction
between reported and subsequently inferred
quantities should remain traceable.


\section{Geometries and Specimens}

The database distinguishes specimen geometry
definitions from the specimen configurations
in which they are used.

The \texttt{geometries} collection contains
reusable geometric configurations associated
with an experimental source.

Each geometry record identifies a geometry
type and specifies the parameter values
required to describe the corresponding
configuration.

Whenever applicable, geometry records
reference the shared CIVITA Geometry Library,
which establishes standardized parameter
names, coordinate systems, dimensional
conventions, and geometric constraints.

A distinction is therefore maintained
between a globally defined geometry type
and a source-specific geometry instance.

The global definition establishes the
scientific meaning of the parametrization,
whereas the geometry instance specifies
the numerical dimensions and other
information associated with the particular
experimental configuration.

The shared geometry definitions are
introduced in Chapter 6 and documented
in Appendix~\ref{app:geometry-library}.

The \texttt{specimens} collection represents
physical specimens or specimen families
defined through associations with the
corresponding material and geometric records.

At a minimum, a specimen configuration
references a material batch and a geometry.
Additional relationships may identify a
reinforcement layout or additive manufacturing
variables when relevant.

The specimen entity therefore provides
the connection between the material
characteristics and the physical
configuration of the experimental object.

For example, several specimens may share
the same material batch while having
different dimensions. Conversely, specimens
with identical geometries may be manufactured
using different material mixtures or
reinforcement configurations.

A specimen record can be associated with
one or more experimental tests. This
allows a common specimen configuration
to be reused in the representation of
different testing procedures or experimental
conditions.

The specimen entity may represent an
individual physical specimen or a family
of specimens sharing a common definition.
The interpretation depends on the
information available from the source
and the manner in which the associated
experimental observations are reported.

This distinction is particularly important
when experimental results are provided
as averages or statistical summaries
of repeated tests.


\section{Reinforcement Materials, Components, and Layouts}

The database represents reinforcement
configurations through three interconnected
collections:
\texttt{reinforcement\_materials},
\texttt{reinforcement\_components}, and
\texttt{reinforcement\_layouts}.

The separation of these entities allows
reinforcement material properties,
individual reinforcing elements, and
complete reinforcement configurations
to be described independently and reused
across related experimental records.

\paragraph{Reinforcement materials.}

The \texttt{reinforcement\_materials}
collection defines the material
characteristics associated with reinforcing
systems.

Depending on the reinforcement type,
the corresponding records may describe
material classification, mechanical
properties, strength characteristics,
elastic moduli, and other relevant
constitutive or physical properties.

A reinforcement material definition may
be referenced by multiple reinforcement
components, allowing consistent use of
material properties across different
reinforcement arrangements.

\paragraph{Reinforcement components.}

The \texttt{reinforcement\_components}
collection represents individual
reinforcement elements or groups of
elements forming part of a complete
reinforcement configuration.

Examples include longitudinal reinforcing
bar layers, stirrup sets, transverse
ties, and externally applied reinforcement
systems.

Each reinforcement component is associated
with a reinforcement material definition
and a reinforcement layout.

A component combines material information
with the geometric and spatial descriptors
necessary to characterize its contribution
to the reinforcement configuration.

Such descriptors may include reinforcement
dimensions, locations, orientations,
spacings, and other parameters relevant
to the particular reinforcement type.

\paragraph{Reinforcement layouts.}

The \texttt{reinforcement\_layouts}
collection represents complete
reinforcement configurations assembled
from individual reinforcement components.

Each layout is associated with a
geometry record and is defined with
respect to the corresponding specimen
coordinate system and geometric
parametrization.

This relationship ensures that the
positions and orientations of
reinforcement components can be
interpreted consistently with
the associated specimen geometry.

A reinforcement layout may include
multiple components representing
different reinforcement systems
or reinforcement regions.

The resulting layout can be referenced
by multiple specimen records sharing
the same geometric and reinforcement
configuration.

The separation among reinforcement
materials, components, and layouts
facilitates the representation of
complex reinforcement systems while
minimizing duplication and supporting
future extensions of the reinforcement
parametrization libraries.


\section{Additive Manufacturing Variables}

The \texttt{printing\_variables} collection
provides a structured representation of
manufacturing parameters associated with
additively manufactured concrete and
cementitious specimens.

Additive manufacturing introduces
process-dependent characteristics that
may significantly influence the
mechanical and physical response of
the resulting material or structural
component.

Relevant information may include
printing speed, extrusion conditions,
material flow, nozzle characteristics,
printing orientation, and other
parameters describing the manufacturing
process.

These variables are represented
separately from the material batch
and specimen geometry, allowing
manufacturing conditions to be
associated explicitly with the
resulting specimen configuration.

A printing variables record may be
referenced by multiple specimens
manufactured using the same
documented process conditions.

The association is optional because
printing variables are not applicable
to conventionally cast or otherwise
manufactured specimens.

The separation of printing variables
from the remaining specimen descriptors
supports systematic investigations
of manufacturing effects and facilitates
the comparison of additively manufactured
specimens with conventional material
and structural configurations.

The corresponding scientific descriptors
and field definitions are documented
in Appendix~\ref{app:experimental-schema}.


\section{Experimental Tests}

The \texttt{tests} collection represents
experimental testing configurations
applied to specimens.

A test record establishes the
relationship between a specimen
configuration and the experimental
procedure used to investigate its
behavior.

Depending on the experimental program,
the test description may include the
testing method, loading configuration,
environmental conditions, observed
physical mechanisms, and other
information required for interpreting
the experimental response.

Examples include uniaxial compression,
direct tension, bending, shear,
torsion, multiaxial loading, and
tests involving coupled mechanical,
thermal, hygral, or chemical processes.

The test entity occupies a central
position in the database hierarchy
because it connects the specimen
configuration to the devices and
supports, loading histories, and
experimental observations associated
with the test.

A single specimen configuration
may be associated with multiple
test records representing different
experimental procedures or conditions.

The current data model also allows
a test definition to represent a
configuration subjected to repeated
experimental execution.

Individual repetition results,
as well as aggregate quantities
obtained from multiple repetitions,
may be represented through
associated data records.

This organization distinguishes
the definition of the experimental
procedure from the observations
obtained through its execution.

Such a distinction is important
for experimental programs involving
multiple nominally identical specimens,
statistical characterization of
material properties, or repeated
testing under comparable conditions.

The standardized representation
of tests facilitates the selection
of experimental evidence associated
with particular loading conditions,
physical mechanisms, or specimen
configurations.

It also supports the reconstruction
of computational boundary-value
problems corresponding to the
experimental investigations.


\section{Devices and Supports}

The \texttt{devices\_and\_supports}
collection represents physical
components associated with
experimental test configurations.

These components may include
supports, loading platens,
actuators, fixtures, sensors,
instrumentation systems, and
other devices relevant to
the application or measurement
of experimental conditions.

Each device or support record
is associated with an experimental
test.

The corresponding description
may include device classification,
geometry, location, orientation,
mechanical characteristics,
or other information relevant
to the experimental procedure.

The explicit representation
of devices and supports is
particularly important when
their characteristics influence
the observed mechanical response.

For example, the dimensions
and mechanical properties of
loading platens, the configuration
of roller supports, and the
arrangement of displacement
measurement devices may affect
the interpretation of an
experimental load--displacement
response.

Device records also establish
references that can be used
by loading histories and
experimental data.

A loading history may identify
the actuator or other device
through which a prescribed
load or displacement is applied.

Similarly, a measurement
record may reference the
instrumentation device
associated with the observed
quantity.

These relationships allow
the database to distinguish
between quantities prescribed
during an experiment and
quantities measured at
specific locations or
through particular devices.

The parametrization of devices
and supports may involve
geometric, mechanical,
and functional descriptors.
When available, standardized
definitions promote consistent
interpretation across
different experimental programs.

Detailed scientific specifications
and supported device types are
addressed through the
corresponding schema and
controlled vocabulary references.


\section{Loading Histories}

The \texttt{loading\_histories}
collection describes the
prescribed loading protocols
or loading time histories
associated with experimental
tests.

A loading history record
is linked to a test and
may additionally reference
a specific device or support
when the prescribed action
is applied through an
identifiable component
of the experimental setup.

Loading histories may
describe force-controlled,
displacement-controlled,
deformation-controlled,
or other experimental
loading procedures.

Depending on the test,
the prescribed history
may involve monotonic,
cyclic, sustained,
transient, or otherwise
time-dependent actions.

The loading history
representation may include
the controlled quantity,
its units, the relevant
independent variable,
and numerical data
describing the prescribed
evolution.

Additional descriptors
may characterize loading
rates, amplitudes,
durations, or other
parameters governing
the experimental procedure.

The distinction between
loading histories and
experimental data is
fundamental.

Loading histories
describe prescribed
actions or experimental
control protocols,
whereas data records
describe observed,
measured, or reported
responses.

For example, a prescribed
actuator displacement
history is represented
as a loading history,
while the corresponding
measured reaction force
is represented as an
experimental observation.

When the actual applied
history differs from
its nominal prescription,
the corresponding measured
information may also
be preserved as
experimental data.

This separation facilitates
the reconstruction of
experimental boundary
conditions and the
comparison of computational
predictions with measured
responses.


\section{Experimental Data}

The \texttt{data} collection
contains the experimental
observations associated
with individual tests.

It constitutes the principal
representation of the
measured and reported
scientific evidence
preserved in the database.

Each data record is
associated with a test
and may optionally
reference a device
or support when the
observation is related
to a specific measurement
component.

Experimental data
may be represented
through different
data structures,
depending on the
nature of the
reported information.

These include scalar
quantities, curves,
time histories,
tables, images,
videos, and
multidimensional
measurement fields.

Scalar observations
may describe quantities
such as peak loads,
strength values,
elastic moduli,
fracture energies,
or other experimentally
determined characteristics.

Curves and time histories
may describe relationships
among applied loads,
displacements, strains,
temperatures, environmental
conditions, or other
measured quantities.

Tabular data may
provide collections
of numerical results
associated with
different experimental
conditions or specimen
configurations.

Images and videos
may document specimen
configurations,
experimental setups,
deformation patterns,
crack development,
failure mechanisms,
or other observations
relevant to the
interpretation of
experimental behavior.

Multidimensional fields
may include spatially
distributed measurements
obtained through
digital image correlation,
thermal imaging,
or other experimental
techniques.

Individual data records
may describe observations
from a particular
experimental repetition
or statistical summaries
of multiple repetitions.

When statistical
information is reported,
the distinction among
individual observations,
mean values, dispersion
measures, and other
aggregate quantities
should remain explicit.

Experimental data
must be interpreted
together with the
associated test,
specimen, material,
and configuration
records.

For example, a
load--displacement
curve cannot be
interpreted adequately
without information
regarding specimen
geometry, experimental
loading conditions,
and the locations
or definitions of
the measured quantities.

This contextual
representation supports
the systematic comparison
of observations obtained
from different experimental
programs and their
subsequent use in
computational model
assessment.

The representation
of numerical data,
associated files,
units, and measurement
conventions is discussed
in Chapter~\ref{chap:experimental-data-standards}.


\section{Identifiers and Referential Integrity}

The CIVITA Experimental
Database adopts standardized
identifiers to establish
unambiguous relationships
among its scientific entities.

Each record is associated
with an identifier
appropriate to its
entity collection.

The current conventions
include identifiers
such as \texttt{SRC0001}
for literature sources,
\texttt{BAT0001}
for material batches,
\texttt{GEO0001}
for geometry records,
\texttt{SPC0001}
for specimens,
\texttt{TST0001}
for tests, and
\texttt{DAT0001}
for experimental
data records.

Additional prefixes
identify printing
variables, reinforcement
materials, reinforcement
components, reinforcement
layouts, devices and
supports, and loading
histories.

These identifiers
establish references
within the corresponding
experimental source
and submission context.

The current hierarchy
requires each dataset
submission to contain
a single literature
source. All associated
records must reference
that source through
the \texttt{source\_id}
field.

Relationships among
records must resolve
within the same source,
ensuring that each
experimental dataset
forms a consistent
and identifiable
collection of
scientific information.

The source association
is distinct from the
scientific parent--child
relationships among
entities.

For example, a
specimen record is
associated with its
literature source
while also referencing
a material batch
and a geometry.
Likewise, a test
record references
both its source
and the specimen
configuration to
which it applies.

This distinction
supports the explicit
representation of
scientific relationships
without relying
exclusively on the
hierarchical nesting
of individual records.

Referential integrity
requires all references
between records to
identify existing
entities within
the applicable
scope.

Required relationships
must be present
and valid, while
optional relationships
may be omitted when
the corresponding
scientific information
is not applicable
or unavailable,
subject to the
requirements of
the relevant schema.

A further distinction
is maintained between
source-specific
identifiers and
globally defined
scientific objects.

For example,
\texttt{GEO0001}
identifies a geometry
record within an
experimental source,
whereas a global
geometry type
identifier, such
as \texttt{GLOB-GEO0001},
identifies a shared
parametrization
definition maintained
by the CIVITA Network.

The geometry record
references the
corresponding global
definition and its
version while retaining
its own source-specific
identifier and numerical
parameter values.

This separation
allows shared
scientific definitions
to evolve independently
of the experimental
records that use them.

During subsequent
database ingestion
and publication,
source-specific
records may be
associated with
additional identifiers
necessary to ensure
persistent and
unambiguous references
across the CIVITA Network.

The corresponding
procedures must
preserve the
relationships
established within
the original dataset
and maintain
traceability to
the submitted
scientific information.

The detailed identifier
formats, field definitions,
and validation
requirements are
provided in
Appendix~\ref{app:experimental-schema},
while the procedures
governing data ingestion
and identifier management
are discussed in
Chapter~\ref{chap:experimental-curation}.




%============================================================
\chapter{Experimental Data Representation and Scientific Standards}
\label{chap:experimental-data-standards}
%============================================================

\section{Metadata and Scientific Descriptors}

\section{Controlled Vocabularies}

\section{Parametrization Libraries}

\section{Units and Coordinate Systems}

\section{Representation of Experimental Measurements}

\subsection{Scalar Quantities}

\subsection{Curves and Time Histories}

\subsection{Tabular Data}

\subsection{Multidimensional and Field Measurements}

\subsection{Images, Videos, and Other Supporting Files}

\section{File Formats and Data Exchange}

\section{Missing Data and Measurement Uncertainty}

\section{Data Quality and Scientific Consistency}


%============================================================
\chapter{Experimental Data Submission}
\label{chap:experimental-submission}
%============================================================

\section{Introduction and General Workflow}

\section{Submission Architecture and Common Data Format}

\section{Manual Data Extraction and Digitization}

\section{Manual Entry of Existing Datasets}

\section{Local AI Data Extraction}

\section{Personal AI and Direct JSON Submission}

\section{Dataset Review and Validation}

\section{Submission and Data Transfer}


%============================================================
\chapter{Experimental Data Curation and Database Ingestion}
\label{chap:experimental-curation}
%============================================================

\section{Introduction and Scope}

\section{Dataset Acceptance and Validation}

\section{Scientific Review and Quality Assurance}

\section{Provenance Verification}

\section{Identifier Management and Referential Integrity}

\section{Database Ingestion and Publication}

\section{Dataset Versioning and Corrections}

\section{Long-Term Preservation and Stewardship}


%============================================================
\chapter{Experimental Database Search and Retrieval}
\label{chap:experimental-search}
%============================================================

\section{Introduction and Scientific Objectives}

\section{Search Application Architecture}

\section{Parameter Discovery and Query Construction}

\section{Boolean Queries and Relationships Among Records}

\section{Search Results and Scientific Interpretation}

\section{Dataset Inspection and Retrieval}

\section{Data Export and Reuse}

\section{Integration with Other CIVITA Components}
